This Vulnerability Assessment and Penetration Testing (VAPT)
engagement was conducted against the OWASP Juice Shop application in
a controlled local testing environment. The assessment followed a
structured methodology consisting of planning and scoping,
reconnaissance, vulnerability scanning and assessment, and manual
verification of identified vulnerabilities.

During the assessment, automated security testing and manual testing
were used to identify and validate security weaknesses in the
application. The identified vulnerabilities were manually verified
and supported with appropriate evidence collected during the
assessment.

\subsection*{Summary of Findings}

A total of eight security vulnerabilities were confirmed during the
assessment. The findings covered multiple security areas, including
access control, authentication, SQL injection, cross-site scripting,
request integrity, and redirect handling.

The confirmed findings are summarized below:

\begin{itemize}
    \item \textbf{BAC -- Broken Access Control:} Privilege escalation
    through the registration functionality.

    \item \textbf{IDOR -- Insecure Direct Object Reference:}
    Unauthorized access to another user's basket by manipulating the
    basket identifier.

    \item \textbf{AUTH -- Weak Password Recovery:} Weaknesses in the
    password recovery mechanism that may affect account security.

    \item \textbf{SQLI-01 -- SQL Injection:} SQL injection affecting
    the login functionality and allowing authentication bypass under
    the tested conditions.

    \item \textbf{SQLI-02 -- SQL Injection:} SQL injection affecting
    the product search functionality and allowing unauthorized
    database interaction under the tested conditions.

    \item \textbf{XSS -- Reflected Cross-Site Scripting:} Execution
    of attacker-controlled JavaScript through the vulnerable search
    functionality.

    \item \textbf{CSRF -- Cross-Site Request Forgery:} Unauthorized
    profile modification through a cross-site request while the
    victim is authenticated.

    \item \textbf{OR -- Open Redirect:} Redirect handling that allows
    an attacker-controlled destination to be used under the tested
    conditions.
\end{itemize}

\subsection*{Risk Summary}

The confirmed findings consisted of three Critical-severity
vulnerabilities, one High-severity vulnerability, and four
Medium-severity vulnerabilities.

The Critical findings affect important application security controls,
including authentication, authorization, and account recovery. The
High and Medium findings demonstrate additional weaknesses affecting
database interaction, client-side security, resource authorization,
request integrity, and redirect handling.

The combination of these findings demonstrates that multiple layers
of application security controls require strengthening.

\subsection*{Recommended Next Steps}

The Critical and High severity findings should be addressed with
priority. Recommended remediation measures include implementing
parameterized database queries, enforcing server-side authorization,
strengthening password recovery, applying appropriate output
encoding, implementing object-level access controls, and protecting
state-changing requests against CSRF.

Following remediation, the application should undergo a dedicated
retesting process. Each confirmed vulnerability should be tested
again using the original proof-of-concept and relevant alternative
test cases to verify that the vulnerability has been effectively
addressed.

\subsection*{Limitations}

This assessment was conducted under the following constraints, which
should be considered when interpreting the findings in this report:

\begin{itemize}
    \item The engagement was performed by a single tester within a
    fixed testing window, rather than a multi-tester team over an
    extended engagement period.

    \item Testing was scoped to the locally hosted OWASP Juice Shop
    instance only; no testing was performed against a production
    deployment, third-party services, or the underlying host
    operating system.

    \item Denial-of-Service (DoS) testing was explicitly excluded
    from the assessment scope and was not performed.

    \item Automated scanning coverage was limited to OWASP ZAP,
    ffuf, and SQLMap; findings not detected by these tools or through
    manual testing within the assessment window may not be reflected
    in this report.

    \item Large-scale exploitation activities, such as bulk data
    extraction, mass account enumeration, or credential-cracking
    against retrieved password hashes, were intentionally not
    performed, consistent with a controlled and non-disruptive
    testing approach.
\end{itemize}

This assessment reflects the security posture of the application at
the time of testing and within the defined scope. It should not be
interpreted as a guarantee that no additional vulnerabilities exist
outside the tested scope or testing methodology.

\subsection*{Final Statement}

This assessment provides a security-focused view of the OWASP Juice
Shop application based on the testing activities performed within
the defined assessment scope. The findings, evidence, risk
assessment, and remediation recommendations documented in this
report provide a basis for improving the application's security
controls.

Successful remediation should be followed by security retesting and
continuous security testing throughout the application development
lifecycle. This helps ensure that previously identified
vulnerabilities remain resolved and that newly introduced security
weaknesses are identified before deployment.